% ===========================================================================
%  figure/tikz/composizione-ensemble.tex
%  Come si compone un ensemble: dagli stream grezzi al PnL finale.
%  I due punti in cui si innestano i plugin sono marcati.
% ===========================================================================
\begin{tikzpicture}[x = 1mm, y = 1mm]

  % ---------- gli stream in ingresso --------------------------------------
  \foreach \i/\y in {1/10, 2/0, 3/-10} {
    \node[archivio, minimum width = 24mm, minimum height = 8mm] at (12,\y)
      (s\i) {stream \i};
  }

  \node[motore, minimum width = 30mm, minimum height = 14mm] at (50,0) (norm)
    {normalizzazione\\{\scriptsize vol trailing, per stream}};

  \node[motore, minimum width = 30mm, minimum height = 14mm] at (92,0) (pesi)
    {pesatura\\{\scriptsize 14 modalità}};

  \node[motore, minimum width = 30mm, minimum height = 14mm] at (92,-30) (over)
    {overlay\\{\scriptsize 7, moltiplicativi}};

  \node[archivio, minimum width = 26mm, minimum height = 9mm] at (44,-30) (out)
    {PnL\\dell'ensemble};

  \foreach \i in {1,2,3} \draw[flusso] (s\i) -- (norm);
  \draw[flusso] (norm) -- (pesi);
  \draw[flusso] (pesi) -- (over);
  \draw[flusso] (over) -- (out);

  % ---------- i due punti d'innesto ---------------------------------------
  \node[astratto, minimum width = 32mm, minimum height = 8mm] at (132,0) (plA)
    {\codT{weights/}};
  \node[astratto, minimum width = 32mm, minimum height = 8mm] at (132,-30) (plB)
    {\codT{overlays/}};
  \draw[dipendenza] (plA) -- (pesi);
  \draw[dipendenza] (plB) -- (over);

  % ---------- le due note che spiegano la differenza ----------------------
  \node[nota, anchor = north, text width = 50mm, align = center] at (92,-9)
    {Distribuisce il capitale.\\Somma \codT{1}.};
  \node[nota, anchor = north, text width = 50mm, align = center] at (92,-39)
    {Scala l'esposizione.\\Prodotto dei moltiplicatori:\\l'ordine non conta.};
  \node[nota, anchor = north, text width = 42mm, align = center] at (50,-9)
    {Ogni stream portato\\alla stessa vol obiettivo.};

\end{tikzpicture}
